
% www.brianrhall.net

% Document settings
\documentclass[11pt]{article}
\usepackage[margin=1in]{geometry}
\usepackage[pdftex]{graphicx}
\usepackage{multirow}
\usepackage{setspace}
\usepackage{listings}
\pagestyle{plain}
\setlength\parindent{0pt}

\begin{document}
	
	 
	\section*{Le transfert thermique}
	
	\subsection*{Les modes de transfert thermique}
	
	Le transfert thermique est un échange d’énergie thermique entre deux systèmes sans travail mécanique.  
	Il existe trois modes de transfert thermique :
	
	\begin{itemize}
		\item \textbf{La conduction} : l’agitation thermique des molécules se propage d’un endroit à un autre. Elle se produit lorsqu’on met en contact un corps chaud et un corps froid (une plaque électrique en est un exemple).
		
		\item \textbf{Le rayonnement} : il s’effectue par émission d’ondes électromagnétiques (par exemple infrarouges). C’est ainsi que l’on reçoit l’énergie thermique du Soleil.
		
		\item \textbf{La convection} : elle se fait grâce aux déplacements de la matière. Elle se produit dans les appareils de chauffage appelés convecteurs. La convection peut être naturelle ou forcée.
	\end{itemize}
	
	Souvent, plusieurs modes de transfert ont lieu simultanément.
	
	\subsection*{Exemple : isolation thermique d'une maison}
	
	Le transfert thermique à travers les murs d’une maison est limité en raison de leur faible conductivité thermique.  
	Pour diminuer la conduction thermique, il faut utiliser des matériaux peu conducteurs (souvent peu denses). Cependant, il faut également limiter la convection entre l’extérieur et l’intérieur.
	
	La solution consiste à utiliser un matériau de faible densité formant une couche épaisse, comme la laine de verre.
	
	\subsection*{Exemple : un thermos}
	
	Dans le cas d’une bouteille isolante, on utilise un récipient à double paroi avec du vide entre les parois.  
	Ce dispositif doit être rigide afin de résister à la pression atmosphérique.
	
	\subsection*{Le flux thermique}
	
	Le flux thermique est la quantité d’énergie thermique transférée par unité de temps.
	
	Considérons le cas d’un mur :
	\begin{itemize}
		\item Le flux thermique est proportionnel à la différence de température.
		\item Il est proportionnel à la surface du mur.
		\item Il est inversement proportionnel à l’épaisseur du mur.
	\end{itemize}
	
	\subsection*{La conductivité thermique}
	
	Loi de Fourier :
	\[
	\Phi = \lambda \times S \times \frac{\Delta T}{e}
	\]
	
	où :
	\begin{itemize}
		\item $\Phi$ est le flux thermique,
		\item $e$ est l’épaisseur du mur,
		\item $\lambda$ est la conductivité thermique du matériau.
	\end{itemize}
	
	La conductivité thermique dépend du matériau et s’exprime en $W \cdot m^{-1} \cdot K^{-1}$.  
	Son inverse est appelé résistance thermique.
	
	\subsection*{La résistance thermique}
	
	\[
	\Phi = \frac{\Delta T}{R_{th}}
	\]
	
	Pour un mur homogène :
	\[
	R_{th} = \frac{e}{\lambda \times S}
	\]
	
	\subsection*{Transfert dans des résistances en série}
	
	Avec deux résistances en série, le flux est le même dans chaque résistance :
	\[
	\Phi_1 = \frac{\Delta T_1}{R_1}, \quad
	\Phi_2 = \frac{\Delta T_2}{R_2}, \quad
	\Phi_1 = \Phi_2 = \Phi
	\]
	
	La différence de température totale est :
	\[
	\Delta T = \Delta T_1 + \Delta T_2
	\]
	
	Donc :
	\[
	R = R_1 + R_2
	\]
	
	\subsection*{Transfert dans des résistances en parallèle (pont thermique)}
	
	Avec deux résistances en parallèle, la différence de température est la même :
	\[
	\Delta T_1 = \Delta T_2 = \Delta T
	\]
	
	Le flux total est :
	\[
	\Phi = \Phi_1 + \Phi_2
	\]
	
	Donc :
	\[
	\Phi = \Delta T \left( \frac{1}{R_1} + \frac{1}{R_2} \right)
	\]
	
	et :
	\[
	R = \frac{R_1 R_2}{R_1 + R_2}
	\]
	
	Par conséquent, si l’une des résistances est faible, la résistance totale l’est également : c’est ce qu’on appelle un pont thermique.
	
\end{document}